\section{Related work}
\label{sec:related_work}

\subsection{Tool verification and execution monitoring}

Runtime verification of tool calls, state snapshots, and execution traces
is standard in software engineering and has been adapted for language-model
agents~\citep{shinn2024,yao2023}. These systems detect execution errors but
do not decide whether verification is cost-effective, do not maintain a
calibrated belief over multiple failure causes, and typically treat
verification as an always-on mechanism rather than a budgeted action. Our
controller treats each verification check as a decision with explicit cost
and expected information gain, enabling comparison with fixed-schedule and
threshold-based alternatives at matched budget.

\subsection{Selective prediction, calibration, and abstention}

Calibrated uncertainty~\citep{guo2017} and selective
prediction~\citep{geifman2017} provide formal frameworks for when to abstain.
These methods typically apply to single-prediction settings with a fixed
cost for abstention. Our setting is sequential: the decision at step $t$
affects what is observed at $t+1$, and the utility depends on the full
trajectory outcome. The controller must therefore balance immediate
verification cost against the risk of downstream irreversibility.

\subsection{Agent recovery and repair}

AgentRx targets critical-step localization and root-cause categories for
diagnosis~\citep{agentrx2026}. Aegis combines taxonomy with intervention
suggestions~\citep{aegis2025}. These systems demonstrate that recovery
actions can be mapped from diagnosed faults but do not separately evaluate
whether the diagnosis-to-recovery mapping is correct, whether the recovery
is independently verified, or whether alternative recovery actions would
have been superior. Our controller records all available actions, their
costs, the selected action, and the selection propensity, enabling offline
counterfactual evaluation where replay validity permits.

\subsection{Process supervision and post-training}

CAPF converts zero-reward failures into verifier-guided repair trajectories
for training~\citep{capf2026}. Afterburner compares SFT, DPO, and GRPO
training under execution feedback on code optimization tasks~\citep{afterburner2025}.
Zelikman et al.~\citep{zelikman2022} introduce process-level reward for
reasoning. These methods share the premise that step-level or process-level
signals carry more information than terminal outcomes. Our contribution is
not another process-supervision algorithm but a controlled comparison of
process-supervised versus outcome-only training where both consume the same
independently adjudicated transitions, both are evaluated at matched data
and compute, and both are tested on unseen fault compositions with a
predeclared clean-task non-inferiority gate.

\subsection{Value of information and budgeted decision-making}

Value-of-information analysis~\citep{howard1966} and budgeted
reinforcement learning provide theoretical frameworks for deciding
when to gather evidence. Our controller adopts a myopic value-of-
information heuristic as one baseline; the learned policy is trained
from the same belief states. The primary evaluation treats the controller
as a policy to be compared at matched cost, not as a claim about optimality.

The key differentiation from prior work is the combination: calibrated
beliefs as fixed inputs (not co-trained), an action space that spans
verification, recovery, and abstention with explicit costs, independent
postcondition adjudication, randomized policy comparison, and a training
evaluation that isolates process information from terminal information
under matched budgets and a regression gate.
